% ECONOMICS_PROBLEM: {"id": "C73-532", "title": "Sharper approximation rates with stochastic action sets", "jel_code": "C73", "source_id": "OP-0153"}
% !TeX program = xelatex
\documentclass[11pt,a4paper]{article}
\usepackage[margin=23mm]{geometry}
\usepackage{fontspec}
\setmainfont{Latin Modern Roman}
\usepackage{amsmath,amssymb,mathtools}
\usepackage{xcolor,enumitem,booktabs,longtable,array,xurl,hyperref}
\definecolor{muted}{HTML}{53636C}
\hypersetup{colorlinks=true,urlcolor=blue,linkcolor=blue}
\setlength{\parindent}{0pt}
\setlength{\parskip}{5pt}
\newcommand{\R}{\mathbb R}
\newcommand{\E}{\mathbb E}
\newcommand{\N}{\mathbb N}
\newcommand{\ind}{\mathbf 1}
\newcommand{\Var}{\operatorname{Var}}
\newcommand{\Cov}{\operatorname{Cov}}
\newcommand{\supp}{\operatorname{supp}}
\DeclareMathOperator*{\argmin}{arg\,min}
\DeclareMathOperator*{\argmax}{arg\,max}
\newcommand{\entrymeta}[3]{{\sffamily\small\color{muted}\textbf{\texttt{#1}}\quad|\quad #2\par #3\par}}
\newcommand{\Problemref}[1]{\texttt{#1}}
\newcommand{\JELref}[2]{\texttt{#2} (JEL \texttt{#1})}
\begin{document}
\section*{Sharper approximation rates with stochastic action sets}
\textbf{Problem ID:} \texttt{C73-532}\quad \textbf{JEL:} \texttt{C73}\par
% BEGIN ECONOMICS STATEMENT
\label{problem:OP-0153}
% GAME_THEORY_PROBLEM: {"local_id": "GT-023", "original_number": 23, "kind": "R", "entry_kind": "statement_only", "published": false, "review_required": true, "category": "game_theory"}
\label{game:GT-023}
{\sffamily\small\color{muted}
{\sffamily\small\color{muted}\textbf{GT-023} | Learning in games\par R | Quantitative specialization of an explicit source question\par}
\par}
\subsection*{Definitions and mathematical statement}
Let $A_1,A_2$ be finite and let $U\in[-1,1]^{A_1\times A_2}$ be a zero-sum payoff matrix. Independently for the two players and independently across rounds, Nature draws a nonempty feasible set $S_i\subseteq A_i$ from a distribution $\rho_i$; player $i$ observes only $S_i$. A feasible policy is a map $\pi_i:S\mapsto\Delta(S)$. Its implemented marginal is
\[
 \mu_i(a)=\sum_{\varnothing\ne S\subseteq A_i}\rho_i(S)\pi_i(a\mid S),
 \qquad \mathcal M_i=\{\mu_i:\pi_i\text{ feasible}\}.
\]
For $\mu_i\in\mathcal M_i$, define strategy-profile regret, a nonnegative saddle-point gap, by
\[
 \operatorname{SPR}(\mu_1,\mu_2)=
 \max_{\nu_1\in\mathcal M_1}\nu_1^{\mathsf T}U\mu_2
 -\min_{\nu_2\in\mathcal M_2}\mu_1^{\mathsf T}U\nu_2.
\]
A positive weight vector $w_i\in\Delta(A_i)$ implements
$\pi_i^{w_i}(a\mid S)=w_i(a)\mathbf1\{a\in S\}/\sum_{b\in S}w_i(b)$.
The source combines its CAT-FTPL method with a stochastic-approximation procedure producing such weights. For failure probability $p\in(0,1)$, Theorem 5.8 gives
\[
 \Pr\!\left(\operatorname{SPR}(\widetilde\mu_1,\widetilde\mu_2)^2
 \leq O\!\left(\frac1{p\sqrt T}+\frac1{p^2T}\right)\right)\geq1-p,
\]
with polynomial action-count factors suppressed. An explicit improvement target is a procedure in the same sampling/feedback model with a bound $b_{d}(T,p)$, $d=|A_1|+|A_2|$, satisfying, uniformly over normalized instances,
\[
 \Pr(\operatorname{SPR}^2\leq b_d(T,p))\geq1-p,
 \qquad b_d(T,p)=o(T^{-1/2})\quad(T\to\infty)
\]
for every fixed $d$ and $p$, while retaining polynomial dependence on $d$.
\subsection*{Short English statement}
Can the stochastic-approximation step give a faster rigorous equilibrium-error bound?
\subsection*{Variants and scope boundaries}
The displayed rate concerns squared SPR. Improving confidence dependence as $p\downarrow0$ is a separate joint-asymptotic target. Private availability, independence, and the full available-action utility feedback are substantive assumptions. Positive weights provide approximate marginal representations; an exact universal representation must not be inferred from an earlier manuscript version.
\subsection*{Source relationship and status}
[S1], version 5, explicitly asks for improved stochastic-approximation bounds immediately after Theorem 5.8. The little-$o$ criterion above is an editorial quantitative target, not a verbatim conjectured optimal rate.
\subsection*{References}
\begingroup\footnotesize\raggedright
\textbf{[S1]} Thomas Schwarz, Jiaru Li, Ryann Sim and Chun Kai Ling, \emph{Computing Equilibria in Games with Stochastic Action Sets}, arXiv:2602.16234v5 (2026). Locators: Definition 3.1 and Assumption 3.3; Theorem 4.5; Section 5.2, Theorem 5.8 and following paragraph, pp. 9--10; Appendix B.13, Algorithm 3. \url{https://arxiv.org/pdf/2602.16234}
\par\endgroup

\subsection*{Common conventions: Source mathematical conventions}

\textbf{Finite games.} For a finite set $A$, $\Delta(A)$ is its probability
simplex. Players choose independent mixed strategies in Nash equilibrium;
correlation is allowed only when explicitly part of the solution concept.
In a game with payoff functions $u_i$ and strategy profile $x$, an additive
$\varepsilon$ Nash equilibrium satisfies
\[
 u_i(x)\geq u_i(x_i',x_{-i})-\varepsilon
 \quad\text{for every player $i$ and every admissible deviation $x_i'$}.
\]
For numerical approximation, payoffs are normalized as locally specified.
An $\varepsilon$ payoff residual is distinct from distance at most
$\varepsilon$ to the set of exact equilibria. Point convergence, convergence
of empirical play, and convergence in distance to an equilibrium set are
also distinct.

\textbf{Computational representations.} Unless stated otherwise, a complexity
question takes rational data with binary encoding; polynomial time is measured
in total bit length. A fixed-accuracy polynomial algorithm is not necessarily
polynomial in $1/\varepsilon$, and neither is necessarily polynomial in
$\log(1/\varepsilon)$. Succinct games and value, demand, or sampling oracles
must declare their interfaces. Exact real solutions require an explicit
representation; an unspecified list of arbitrary real numbers is not a
finite computational output. A claimed algorithmic barrier is meaningful
only for its specified model and reductions.

\textbf{Dynamic games.} Histories, information, observation, and allowable
randomization are part of the model. A stationary strategy depends only on
the current state; a positional strategy in a deterministic graph game
chooses one action at each controlled vertex. Nash equilibrium at an initial
state and equilibrium after every history are different requirements.
An expectation of a pathwise $\liminf$ payoff, a $\liminf$ of expected averages,
a limiting expected average, and a uniform finite-horizon equilibrium are
different criteria. None is silently substituted for another.

\textbf{Allocation and mechanisms.} A complete allocation assigns every item
exactly once unless otherwise stated. Zero-valued goods, negative values,
capacity constraints, feasibility, lotteries, and copies of goods affect
fairness definitions and are specified locally. Dominant-strategy incentive
compatibility, Bayesian incentive compatibility, universal truthfulness,
and truthfulness in expectation impose different inequalities. Whenever
an objective ratio has zero denominator, its convention is stated or those
instances are excluded.

\textbf{Infinite-dimensional models.} Expectations, integrals, stochastic
equations, and optimization objectives require their local regularity,
integrability, filtrations, and admissible domains. A backward--forward
mean-field system is a boundary-value problem; it is not an initial-value
semiflow without an additional construction. Long-time, large-population,
and vanishing-noise limits require an order of limits and a convergence
topology. If a minimum need not be attained, the objective is its infimum
or supremum and the approximation requirement is stated separately.
% END ECONOMICS STATEMENT
\end{document}
